\chapter{Complexity — BPL In P}
%%% generated by the blueprint pipeline from TCSlib.Complexity.Randomized.BPLInP
\section{Overview}

Declarations of the Lean module \texttt{TCSlib.Complexity.Randomized.BPLInP}.

\section{Declarations}

\begin{definition}[Context word]
\lean{Randomized.BPLInP.ctxW}
\label{Randomized.BPLInP.ctxW}
\leanok
\uses{Randomized.BPLInP.codesW, Randomized.BPLInP.radius, Turing.pairEncode}
For a nondeterministic finite-state machine $N$ and natural numbers $c$, $B$, and $e$,
the context word associated with a Boolean list $x$ is the pair encoding of $x$ together
with the pair encoding of the radius determined by $c$ and $x$ and the $W$-codes
determined by $B$, $e$, and $x$.
\end{definition}

\begin{definition}[One-step transition test]
\lean{Randomized.BPLInP.stepTest}
\label{Randomized.BPLInP.stepTest}
\leanok
\uses{Complexity.Codec.stepLang, Complexity.pairFstD, Complexity.pairSndD, Randomized.BPLInP.fixCoin, Turing.MultiTapeTM.indicator, Turing.pairEncode}
Let $N$ be a finite nondeterministic Turing machine with Boolean-valued tapes, let
$c,B,e$ be natural numbers, let $b$ be a Boolean, and let $z$ be a Boolean list. The
test returns true exactly when the nested pair encoding obtained from the indicated
first and second projections of $z$ belongs to the one-step language of $N$ with its
coin fixed to $b$.
\end{definition}

\begin{definition}[Encoded transition entry]
\lean{Randomized.BPLInP.entryW}
\label{Randomized.BPLInP.entryW}
\leanok
\uses{Complexity.BinArith.add, Complexity.pairSndD, Randomized.BPLInP.stepTest, Turing.pairEncode}
Let $N$ be a finite nondeterministic Turing machine with Boolean-valued tapes, let
$c,B,e$ be natural numbers, and let $z$ be a Boolean list. Form a pair whose first
component is the second projection of the second projection of $z$ when the one-step
test for coin value false accepts $z$, and is empty otherwise; form its second component
in the same way using coin value true. Return the binary addition of the encoding of
this pair.
\end{definition}

\begin{definition}[Encoded transition row]
\lean{Randomized.BPLInP.rowW}
\label{Randomized.BPLInP.rowW}
\leanok
\uses{Complexity.BinArith.zipWithW, Complexity.pairFstD, Complexity.pairSndD, Randomized.BPLInP.entryW, Randomized.BPLInP.sumW, Turing.pairEncode}
Let $N$ be a finite nondeterministic Turing machine with Boolean-valued tapes, let
$c,B,e$ be natural numbers, and let $y$ be a Boolean list. Encode two pairs formed from
the indicated first and second projections of $y$; zip the resulting list with the list
obtained by encoding the first pair, apply the encoded transition-entry operation to
corresponding entries, and sum the resulting list.
\end{definition}

\begin{definition}[Dynamic-programming transition step]
\lean{Randomized.BPLInP.dpStep}
\label{Randomized.BPLInP.dpStep}
\leanok
\uses{Complexity.BinArith.mapW, Complexity.pairFstD, Complexity.pairSndD, Randomized.BPLInP.rowW, Turing.pairEncode}
Let $N$ be a finite nondeterministic Turing machine with Boolean-valued tapes, let
$c,B,e$ be natural numbers, and let $S$ be a Boolean list. Return an encoding whose
first component is the first projection of $S$ and whose second component is the result
of applying the row operation to the list obtained by encoding $S$ together with the
specified nested projections of $S$.
\end{definition}

\begin{definition}[Initial acceptance entry]
\lean{Randomized.BPLInP.initEntry}
\label{Randomized.BPLInP.initEntry}
\leanok
\uses{Complexity.Codec.acceptLang, Complexity.pairFstD, Complexity.pairSndD, Turing.MultiTapeTM.indicator, Turing.pairEncode}
Let $N$ be a finite nondeterministic Turing machine with Boolean-valued tapes, let
$c,B,e$ be natural numbers, and let $y$ be a Boolean list. Return the singleton list
containing true if the nested pair encoding formed from the indicated projections of $y$
belongs to the acceptance language of $N$, and return the empty list otherwise.
\end{definition}

\begin{definition}[Dynamic-programming initialization]
\lean{Randomized.BPLInP.dpInit}
\label{Randomized.BPLInP.dpInit}
\leanok
\uses{Complexity.BinArith.mapW, Randomized.BPLInP.codesW, Randomized.BPLInP.ctxW, Randomized.BPLInP.initEntry, Turing.pairEncode}
Let $N$ be a finite nondeterministic Turing machine with Boolean-valued tapes, let
$c,B,e$ be natural numbers, and let $x$ be a Boolean list. Return an encoding whose
first component is the context associated to $c,B,e,x$ and whose second component is
obtained by applying the initial-entry operation to the list formed from that context
and the codes associated to $B,e,x$.
\end{definition}

\begin{definition}[Dynamic-programming run]
\lean{Randomized.BPLInP.dpRun}
\label{Randomized.BPLInP.dpRun}
\leanok
\uses{Randomized.BPLInP.dpInit, Randomized.BPLInP.dpStep}
Let $N$ be a finite nondeterministic Turing machine with Boolean-valued tapes, let
$c,B,e$ be natural numbers, and let $x$ be a Boolean list. Starting from the
dynamic-programming initialization for $x$, apply the dynamic-programming transition
step exactly $B(|x|+1)^e$ times.
\end{definition}

\begin{definition}[Initial witness selection]
\lean{Randomized.BPLInP.pickW}
\label{Randomized.BPLInP.pickW}
\leanok
\uses{Complexity.Codec.initLang, Complexity.pairFstD, Complexity.pairSndD, Turing.MultiTapeTM.indicator, Turing.pairEncode}
Let $N$ be a finite nondeterministic Turing machine with Boolean-valued tapes, let
$c,B,e$ be natural numbers, and let $z$ be a Boolean list. If the nested pair encoding
formed from the indicated projections of $z$ belongs to the initial language of $N$,
return the second projection of the second projection of $z$; otherwise return the empty
list.
\end{definition}

\begin{definition}[Weighted acceptance string]
\lean{Randomized.BPLInP.accW}
\label{Randomized.BPLInP.accW}
\leanok
\uses{Complexity.BinArith.zipWithW, Complexity.pairFstD, Complexity.pairSndD, Randomized.BPLInP.dpRun, Randomized.BPLInP.pickW, Randomized.BPLInP.sumW, Turing.pairEncode}
Let $N$ be a finite nondeterministic Turing machine over Boolean symbols, let $c,B,e$ be
natural numbers, and let $x$ be a Boolean string. Define the weighted acceptance string
for $x$ by applying the machine’s choice weights to the paired data extracted from the
bounded dynamic-programming run of $N$ on $x$, and summing the resulting weighted
strings.
\end{definition}

\begin{definition}[Bounded probabilistic test]
\lean{Randomized.BPLInP.bplTest}
\label{Randomized.BPLInP.bplTest}
\leanok
\uses{Complexity.BinArith.lt, Randomized.BPLInP.accW, Randomized.BPLInP.pow2W, Turing.pairEncode}
Let $N$ be a finite nondeterministic Turing machine over Boolean symbols, let $c,B,e$ be
natural numbers, and let $x$ be a Boolean string. The test returns whether the binary
encoding of $2^{Be}$ paired with the weighted acceptance string, preceded by a zero bit,
is less than the specified binary encoding comparison bound.
\end{definition}

\begin{theorem}[Polynomial-Time Computability of the Step Test]
\lean{Randomized.BPLInP.polyTimeComputable_stepTest}
\label{Randomized.BPLInP.polyTimeComputable_stepTest}
\leanok
\uses{Complexity.Codec.stepLang_mem, Complexity.PolyTimeComputable, Complexity.PolyTimeComputable.pairEncode, Complexity.polyTimeComputable_pairFstD, Complexity.polyTimeComputable_pairSndD, Randomized.BPLInP.fixCoin, Randomized.BPLInP.polyTimeComputable_indicator_of_LOGSPACE, Randomized.BPLInP.stepTest}
\difficulty{4}
For every finite-state nondeterministic Turing machine $N$ over the binary alphabet,
natural numbers $c$, $B$, and $e$, and bit $b$, the function that maps a bit string $z$
to the one-bit string containing the step test for $N$, coin choice $b$, and input $z$
is polynomial-time computable.
\end{theorem}

\begin{theorem}[Polynomial-Time Computability of a Dynamic-Programming Step]
\lean{Randomized.BPLInP.polyTimeComputable_dpStep}
\label{Randomized.BPLInP.polyTimeComputable_dpStep}
\leanok
\uses{Complexity.BinArith.polyTimeComputable_add, Complexity.BinArith.polyTimeComputable_mapW, Complexity.BinArith.polyTimeComputable_zipWithW, Complexity.PolyTimeComputable, Complexity.PolyTimeComputable.pairEncode, Complexity.pairSndD, Complexity.polyTimeComputable_const, Complexity.polyTimeComputable_id, Complexity.polyTimeComputable_ite, Complexity.polyTimeComputable_pairFstD, Complexity.polyTimeComputable_pairSndD, Randomized.BPLInP.dpStep, Randomized.BPLInP.entryW, Randomized.BPLInP.polyTimeComputable_stepTest, Randomized.BPLInP.polyTimeComputable_sumW, Randomized.BPLInP.rowW}
\difficulty{5}
For every finite-state nondeterministic Turing machine $N$ over the binary alphabet and
all natural numbers $c$, $B$, and $e$, the function that performs one
dynamic-programming step on an encoded bit string is polynomial-time computable.
\end{theorem}

\begin{theorem}[Polynomial-time computability of initialization]
\lean{Randomized.BPLInP.polyTimeComputable_dpInit}
\label{Randomized.BPLInP.polyTimeComputable_dpInit}
\leanok
\uses{Complexity.BinArith.polyTimeComputable_mapW, Complexity.Codec.acceptLang_mem, Complexity.PolyTimeComputable, Complexity.PolyTimeComputable.pairEncode, Complexity.polyTimeComputable_const, Complexity.polyTimeComputable_id, Complexity.polyTimeComputable_ite, Complexity.polyTimeComputable_pairFstD, Complexity.polyTimeComputable_pairSndD, Randomized.BPLInP.ctxW, Randomized.BPLInP.dpInit, Randomized.BPLInP.initEntry, Randomized.BPLInP.polyTimeComputable_codesW, Randomized.BPLInP.polyTimeComputable_indicator_of_LOGSPACE, Randomized.BPLInP.polyTimeComputable_radius}
\difficulty{4}
For every finite nondeterministic Turing machine $N$ over bits and natural numbers
$c,B,e$, the function mapping a bit string to its initialized dynamic-programming
encoding is polynomial-time computable.
\end{theorem}

\begin{theorem}[Polynomial-time test from a computable run]
\lean{Randomized.BPLInP.polyTimeComputable_bplTest_of_run}
\label{Randomized.BPLInP.polyTimeComputable_bplTest_of_run}
\leanok
\uses{Complexity.BinArith.polyTimeComputable_lt, Complexity.BinArith.polyTimeComputable_zipWithW, Complexity.Codec.initLang_mem, Complexity.PolyTimeComputable, Complexity.PolyTimeComputable.pairEncode, Complexity.polyTimeComputable_const, Complexity.polyTimeComputable_ite, Complexity.polyTimeComputable_pairFstD, Complexity.polyTimeComputable_pairSndD, Complexity.polyTimeComputable_prepend, Randomized.BPLInP.accW, Randomized.BPLInP.bplTest, Randomized.BPLInP.dpRun, Randomized.BPLInP.pickW, Randomized.BPLInP.polyTimeComputable_indicator_of_LOGSPACE, Randomized.BPLInP.polyTimeComputable_pow2W, Randomized.BPLInP.polyTimeComputable_sumW}
\difficulty{4}
Let $N$ be a finite nondeterministic Turing machine over bits, and let $c,B,e$ be
natural numbers. If the map sending each bit string $x$ to the output of the associated
dynamic-programming run is polynomial-time computable, then the function sending $x$ to
the one-bit output containing the value of the test on $x$ is polynomial-time
computable.
\end{theorem}

\begin{theorem}[Sum over a range]
\lean{Randomized.BPLInP.sum_map_range}
\label{Randomized.BPLInP.sum_map_range}
\leanok
\difficulty{4}
For every function $f:\mathbb{N}\to\mathbb{N}$ and every $V\in\mathbb{N}$, the sum of
the list $f(0),f(1),\ldots,f(V-1)$ equals the finite sum $\sum_{i=0}^{V-1}f(i)$.
\end{theorem}

\begin{theorem}[Mapping lists through binary expansions]
\lean{Randomized.BPLInP.map_bits}
\label{Randomized.BPLInP.map_bits}
\leanok
\difficulty{1}
For every function $f:\mathbb{N}\to\mathbb{N}$ and every finite list $l$ of natural
numbers, mapping $l$ by $a\mapsto$ the binary expansion of $f(a)$ gives the same list of
bit strings as first mapping $l$ by $f$ and then taking each binary expansion.
\end{theorem}

\begin{definition}[Context encoding]
\lean{Randomized.BPLInP.ctxOf}
\label{Randomized.BPLInP.ctxOf}
\leanok
\uses{Complexity.BinArith.encWords, Turing.pairEncode}
Let $x$ be a Boolean string and let $s,V$ be natural numbers. Define the context string
by pairing $x$ with a pair consisting of the binary representation of $s$ and the
encoding of the binary representations of all integers from $0$ through $V-1$.
\end{definition}

\begin{definition}[Encoded configuration vector]
\lean{Randomized.BPLInP.vecW}
\label{Randomized.BPLInP.vecW}
\leanok
\uses{Complexity.BinArith.encWords, Randomized.BPLInP.dpVal}
Let $N$ be a finite nondeterministic Turing machine over Boolean symbols, let $x$ be a
Boolean string, and let $s,V,t$ be natural numbers. Define the vector string by
encoding, in order for each $a<V$, the binary representation of the value of the bounded
dynamic-programming computation for $N$ on $x$ at parameters $s,V,t,a$.
\end{definition}

\begin{theorem}[Equality of the configuration contexts]
\lean{Randomized.BPLInP.ctxW_eq}
\label{Randomized.BPLInP.ctxW_eq}
\leanok
\uses{Randomized.BPLInP.codesW_eq, Randomized.BPLInP.ctxOf, Randomized.BPLInP.ctxW, Randomized.BPLInP.radius}
\difficulty{3}
Let $x$ be a finite bit string and let $c,B,e$ be natural numbers. The context encoding
of $x$ with radius parameter $c$ and code-iteration parameters $B,e$ equals the context
encoding of $x$ with radius $c((\operatorname{bits}( |\operatorname{bits}(|x|)|)+1)$ and
iteration bound $B(|x|+1)^e$, where the radius parameter is represented by its binary
expansion and the code sequence is iterated for the stated bound.
\end{theorem}

\begin{theorem}[Encoded transition-entry value]
\lean{Randomized.BPLInP.entryW_eq}
\label{Randomized.BPLInP.entryW_eq}
\leanok
\uses{Complexity.BinArith.add, Complexity.Codec.StepRel, Complexity.LogProg.bitsVal, Complexity.LogProg.bitsVal_bits, Complexity.pairFstD_pairEncode, Complexity.pairSndD_pairEncode, Randomized.BPLInP.entryW, Randomized.BPLInP.fixCoin, Randomized.BPLInP.stepTest, Turing.MultiTapeTM.indicator, Turing.pairEncode}
\difficulty{4}
Let $N$ be a finite-state nondeterministic Turing machine over the Boolean alphabet, let
$x$ be a finite bit string, and let $s,a,a',v$ be natural numbers. The encoded entry for
the pair of configurations with codes $a$ and $a'$ is the binary representation of the
sum of two terms: each term is $v$ if the corresponding transition relation, with the
choice bit fixed to $0$ or $1$, holds from code $a$ to code $a'$ on input $x$ within
space bound $s$, and is $0$ otherwise.
\end{theorem}

\begin{theorem}[Initial dynamic-programming encoding]
\lean{Randomized.BPLInP.dpInit_eq}
\label{Randomized.BPLInP.dpInit_eq}
\leanok
\uses{Complexity.BinArith.decWords_encWords, Complexity.BinArith.encWords, Complexity.BinArith.mapW_pairEncode, Complexity.pairFstD_pairEncode, Complexity.pairSndD_pairEncode, Randomized.BPLInP.codesW_eq, Randomized.BPLInP.ctxW, Randomized.BPLInP.dpInit, Randomized.BPLInP.dpVal, Randomized.BPLInP.initEntry, Randomized.BPLInP.radius, Randomized.BPLInP.vecW, Turing.MultiTapeTM.indicator, Turing.pairEncode}
\difficulty{4}
Let $N$ be a finite-state nondeterministic Turing machine over the Boolean alphabet, let
$c,B,e$ be natural numbers, and let $x$ be a finite bit string. The initial
dynamic-programming encoding is the pair of the context encoding of $x$ and the word
encoding, for each $a=0,\ldots,V-1$, of the base value that is $1$ exactly when the
context encoding of $x$ with space bound $s$ and code $\operatorname{bits}(a)$ belongs
to the machine's accepting-configuration language, where
$s=c(|\operatorname{bits}(|x|)|+1)$ and $V=B(|x|+1)^e$.
\end{theorem}

\begin{theorem}[One-step dynamic-programming update]
\lean{Randomized.BPLInP.dpStep_eq}
\label{Randomized.BPLInP.dpStep_eq}
\leanok
\uses{Complexity.BinArith.decWords_encWords, Complexity.BinArith.encWords, Complexity.BinArith.mapW_pairEncode, Complexity.BinArith.zipWithW_pairEncode, Complexity.pairFstD_pairEncode, Complexity.pairSndD_pairEncode, Randomized.BPLInP.ctxOf, Randomized.BPLInP.dpStep, Randomized.BPLInP.dpVal, Randomized.BPLInP.entryW_eq, Randomized.BPLInP.map_bits, Randomized.BPLInP.rowW, Randomized.BPLInP.sumW_encWords, Randomized.BPLInP.sum_map_range, Randomized.BPLInP.vecW, Turing.pairEncode}
\difficulty{5}
Let $N$ be a finite-state nondeterministic Turing machine over the Boolean alphabet, let
$x$ be a finite bit string, and let $s,V,t$ be natural numbers. Applying one encoded
dynamic-programming update to the pair consisting of the context encoding of $x$ with
parameters $s,V$ and the encoded vector of values at depth $t$ yields the same pair with
that context and the encoded vector of values at depth $t+1$.
\end{theorem}

\begin{theorem}[Iterated Dynamic-Programming Update]
\lean{Randomized.BPLInP.iterate_dpStep}
\label{Randomized.BPLInP.iterate_dpStep}
\leanok
\uses{Randomized.BPLInP.ctxOf, Randomized.BPLInP.dpStep, Randomized.BPLInP.dpStep_eq, Randomized.BPLInP.vecW, Turing.pairEncode}
\difficulty{4}
Let $N$ be a finite nondeterministic Turing machine with Boolean choices, let $x$ be a
bit string, and let $s,V,t\in\mathbb N$. Applying the dynamic-programming update $t$
times to the pairing of the context for $x,s,V$ and the initial encoded value vector
gives the pairing of that same context with the encoded vector of dynamic-programming
values at time $t$.
\end{theorem}

\begin{theorem}[Iterated Dynamic-Programming Initialization]
\lean{Randomized.BPLInP.iterate_dpInit}
\label{Randomized.BPLInP.iterate_dpInit}
\leanok
\uses{Randomized.BPLInP.ctxW, Randomized.BPLInP.ctxW_eq, Randomized.BPLInP.dpInit, Randomized.BPLInP.dpInit_eq, Randomized.BPLInP.dpStep, Randomized.BPLInP.iterate_dpStep, Randomized.BPLInP.vecW, Turing.pairEncode}
\difficulty{3}
Let $N$ be a finite nondeterministic Turing machine with Boolean choices, let
$c,B,e,i\in\mathbb N$, and let $x$ be a bit string. Put
$s=c\bigl(|\mathrm{bits}(|x|)|+1\bigr)$ and $V=B(|x|+1)^e$. After $i$ iterations of the
dynamic-programming update, its encoding initialized from $N,c,B,e,x$ is the pairing of
the context encoding for these parameters and the encoded vector whose entry at each
index $a<V$ is the dynamic-programming value at time $i$ and index $a$.
\end{theorem}

\begin{theorem}[Bounded-space test and path-count threshold]
\lean{Randomized.BPLInP.bplTest_eq}
\label{Randomized.BPLInP.bplTest_eq}
\leanok
\uses{Complexity.BinArith.decWords_encWords, Complexity.BinArith.encWords, Complexity.BinArith.lt, Complexity.BinArith.zipWithW_pairEncode, Complexity.Codec.cfgCode, Complexity.Codec.cfgCode_lt, Complexity.Codec.codeK, Complexity.Codec.mem_initLang, Complexity.LogProg.bitsVal, Complexity.LogProg.bitsVal_bits, Complexity.logSpace, Complexity.pairFstD_pairEncode, Complexity.pairSndD_pairEncode, Randomized.BPLInP.accW, Randomized.BPLInP.bitsVal_pow2W, Randomized.BPLInP.bplTest, Randomized.BPLInP.ctxOf, Randomized.BPLInP.ctxW_eq, Randomized.BPLInP.dpRun, Randomized.BPLInP.dpVal, Randomized.BPLInP.iterate_dpInit, Randomized.BPLInP.map_bits, Randomized.BPLInP.pickW, Randomized.BPLInP.sumW_encWords, Randomized.BPLInP.sum_map_range, Randomized.BPLInP.vecW, Turing.MultiTapeTM.indicator, Turing.NDTM.initCfg, Turing.pairEncode}
\difficulty{5}
Let $N$ be a finite nondeterministic Turing machine over Boolean symbols, let $x$ be a
finite Boolean string, and let $c,B,e$ be natural numbers. If
\[
2^{(2W+2)\left(c\bigl(|\operatorname{bits}(|x|)|+1\bigr)+\operatorname{logSpace}(|x|)+1\right)}\le
B(|x|+1)^e,
\]
then the Boolean test on $x$ is true exactly when
\[
2^{B(|x|+1)^e}<2D_{B(|x|+1)^e}\!\left(C_0\right),
\]
where $W$ is the number of bits used to encode a machine configuration cell, $D_t(a)$ is
the recursively defined count of accepting computation paths of length at most $t$
starting from configuration code $a$ and retaining only successor codes below
$B(|x|+1)^e$, and $C_0$ is the code of the initial configuration with space window
$c(|\operatorname{bits}(|x|)|+1)$.
\end{theorem}

\begin{theorem}[Encoded dynamic program length bound]
\lean{Randomized.BPLInP.length_vecW_le}
\label{Randomized.BPLInP.length_vecW_le}
\leanok
\uses{Complexity.BinArith.length_encWords_le, Complexity.LogProg.length_bits_le, Randomized.BPLInP.dpVal, Randomized.BPLInP.dpVal_le, Randomized.BPLInP.vecW}
\difficulty{4}
For every finite nondeterministic Turing machine $N$, bit string $x$, and natural
numbers $s,V,t$, the word produced by the dynamic-programming encoding has length at
most $V\bigl(2(2Vt+1)+2\bigr)$.
\end{theorem}

\begin{theorem}[Length Bound for Context Encoding]
\lean{Randomized.BPLInP.length_ctxOf_le}
\label{Randomized.BPLInP.length_ctxOf_le}
\leanok
\uses{Randomized.BPLInP.ctxOf, Randomized.BPLInP.length_encWords_range_le, Turing.length_bits_le_self, Turing.length_pairEncode}
\difficulty{3}
For every pair of natural numbers $s,V$ and every bit string $x$, the context encoding
of $x,s,V$ has length at most $2|x|+2s+4+V(2V+2)$.
\end{theorem}

\begin{theorem}[Polynomial-time computability of the run]
\lean{Randomized.BPLInP.polyTimeComputable_dpRun}
\label{Randomized.BPLInP.polyTimeComputable_dpRun}
\leanok
\uses{Complexity.BinArith.polyTimeComputable_iterate, Complexity.PolyTimeComputable, Randomized.BPLInP.ctxOf, Randomized.BPLInP.ctxW_eq, Randomized.BPLInP.dpRun, Randomized.BPLInP.iterate_dpInit, Randomized.BPLInP.length_ctxOf_le, Randomized.BPLInP.length_vecW_le, Randomized.BPLInP.polyTimeComputable_dpInit, Randomized.BPLInP.polyTimeComputable_dpStep, Randomized.BPLInP.vecW, Turing.length_bits_le_self, Turing.length_pairEncode}
\difficulty{6}
For every finite nondeterministic Turing machine $N$ over bits and natural numbers
$c,B,e$, the function mapping a bit string to the result of the associated
dynamic-programming run is polynomial-time computable.
\end{theorem}

\begin{theorem}[Polynomial-time computability of the test]
\lean{Randomized.BPLInP.polyTimeComputable_bplTest}
\label{Randomized.BPLInP.polyTimeComputable_bplTest}
\leanok
\uses{Complexity.PolyTimeComputable, Randomized.BPLInP.bplTest, Randomized.BPLInP.polyTimeComputable_bplTest_of_run, Randomized.BPLInP.polyTimeComputable_dpRun}
\difficulty{1}
For every finite nondeterministic Turing machine $N$ over bits and natural numbers
$c,B,e$, the function sending a bit string $x$ to the one-bit output containing the
value of the test on $x$ is polynomial-time computable.
\end{theorem}

\begin{theorem}[Exponential bound for machine codes]
\lean{Randomized.BPLInP.two_pow_codes_le}
\label{Randomized.BPLInP.two_pow_codes_le}
\leanok
\uses{Complexity.Codec.codeK, Complexity.logSpace}
\difficulty{5}
Let $N$ be a nondeterministic Turing machine over the Boolean alphabet, with $k$ work
tapes and finite state set $Q$. Set
$K=2\bigl(\lvert\mathrm{Option}(Q)\times\{0,1,2\}\rvert+3k+1\bigr)+2$, where
$\lvert\cdot\rvert$ denotes cardinality. For every $c,n\in\mathbb{N}$, if $b(n)$ is the
length of the binary expansion of $n$ and $\ell(n)=\lfloor\log_2 n\rfloor+1$ for $n>0$,
with $\ell(0)=1$, then
\[
2^{K\left(c(b(n)+1)+\ell(n)+1\right)}\leq 8^{K(c+1)}(n+1)^{K(c+1)}.
\]
\end{theorem}

\begin{theorem}[BPL is contained in P]
\lean{Randomized.BPL_subset_P}
\label{Randomized.BPL_subset_P}
\leanok
\uses{Complexity.Codec.cfgCode, Complexity.Codec.codeK, Complexity.P, Complexity.logSpace, Complexity.mem_P_of_test, Randomized.BPL, Randomized.BPLInP.bplTest, Randomized.BPLInP.bplTest_eq, Randomized.BPLInP.dpVal, Randomized.BPLInP.dpVal_init, Randomized.BPLInP.haltsWithin_of_codes, Randomized.BPLInP.polyTimeComputable_bplTest, Randomized.BPLInP.two_pow_codes_le, Randomized.randProb_le_one, Randomized.randProb_or_disjoint, Turing.FinNDTM.outputProb, Turing.FinNDTM.outputProb_eq_of_haltsWithin, Turing.MultiTapeTM.indicator, Turing.NDTM.HaltsWithin, Turing.NDTM.initCfg}
\difficulty{6}
Every Boolean language decidable in bounded-error probabilistic logarithmic space is
decidable in deterministic polynomial time.
\end{theorem}
